% Sección independiente; NO modifica Overleaf. Requiere amsmath, booktabs, graphicx, biblatex.
\section{Ingeniería de detalles del prototipo manual}
\subsection{Modelo de compactación y presión}
Se adopta accionamiento manual por manivela de radio nominal y CAD $165\,\mathrm{mm}$.
El motor de $0,18\,\mathrm{kW}$ se conserva como antecedente y no como accionamiento nominal.
La decisión sigue Ingeniería de Detalles, sección1.1\parencite{equipo_detalles_prototipo_2026}.
La ecuación de Holm reproducida por Stelte et al.\ \parencite{stelte_presion_2011} permite estimar
\begin{equation}P=\frac{P_{N0}}{\nu_{LR}}\left[\exp\left(4\mu_w\nu_{LR}\frac{L}{D}\right)-1\right].\end{equation}
Los coeficientes $P_{N0}=0,281\,\mathrm{MPa}$, $\nu_{LR}=0,203$ y $\mu_w=0,715$
se recuperan de EP2, Anexo A.2.2, ecuación4\parencite{equipo_entrega2_2026}.
EP2 los adopta para rastrojo de maíz al20\,\% de humedad; su transferencia al aserrín requiere validación.
Con $L=10\,\mathrm{mm}$ y $D=6\,\mathrm{mm}$ se obtiene $P=2,25867\,\mathrm{MPa}$.
La estimación no constituye medición ni calibración del aserrín.
\subsection{Solicitaciones y accionamiento manual}
La matriz CAD tiene $102$ canales y dos rodillos $52\times30\,\mathrm{mm}$,
a radios $27,838$ y $27,662\,\mathrm{mm}$, con separación geométrica $0,5\,\mathrm{mm}$.
Se conservan tres capas supuestas como alternativas, sin seleccionar una condición nominal física.
El reparto parcial de área y los torques se obtienen del núcleo EP2 con $1440$ estados correlacionados.
La fuerza manual ideal supone aplicación tangencial y $F_h=T/0,165$; no incluye pérdidas ni atasco.
\begin{table}[ht]\centering
\begin{tabular}{rrrrr}\toprule
Capa [mm]&$F_{\max}$ [N]&$\bar F$ [N]&$T_{\max}$ [N m]&$F_{h,\max}$ [N]\\\midrule
1 & 389.30 & 350.69 & 1.1566 & 7.01\\
2 & 641.44 & 603.69 & 1.9335 & 11.72\\
4 & 927.99 & 910.02 & 2.8068 & 17.01\\
\bottomrule\end{tabular}
\caption{Alternativas condicionadas con presión teórica y $\mu_{\mathrm{ef}}=0,1$.}
\end{table}
El cálculo anterior $3,255/0,165=19,727\,\mathrm{N}$ es aritméticamente consistente;
sus cargas se conservan como referencia documental, no como simulación del CAD actual.
Las remisiones de Detalles no coinciden con la numeración de EP2; los anexos A.2.1/A.2.2 disponibles
definen el predimensionamiento comercial y no proporcionan un historial angular del prototipo.
\subsection{Verificación de componentes}
El cribado Euler--Bernoulli del eje transversal considera apoyos articulados rígidos hipotéticos
en taladros de carcasa, conservando los centros CAD de los cuatro rodamientos.
La sensibilidad compara el cilindro liso de $6\,\mathrm{mm}$ y el núcleo de referencia
M6 por 1 de $4,773\,\mathrm{mm}$\parencite{tr_rosca_m6_2026}.
No se conocen grado, raíz real ni rigidez del montaje. La deformación calculada es relevante
respecto de la holgura, por lo que el contacto desacoplado exige revisión física.
Los asientos presentan diferencias diametrales de $2\,\mathrm{mm}$ en ambas interfaces
eje--608 y 608--rodillo. La comparación EZO\parencite{ezo_608zz_2026} no certifica los rodamientos adquiridos.
El árbol se comprueba solo como sección lisa axial--torsional condicionada;
su flexión y apoyos no están resueltos. UCF201-8 y UCFL201-8 permanecen sin selección.
\subsection{Limitaciones y validación pendiente}
No se aprueba resistencia, fatiga, unión roscada, carcasa ni soporte inferior.
Detalles menciona ABS y EP2 identifica PLA: debe confirmarse el filamento físico.
No se ejecuta FEM de la matriz sin justificar contacto y retención; las figuras CAD no son FEM.
El protocolo de aserrín prescribe humedad de $12$--$15\,\%$ en base húmeda.
Para $500\,\mathrm{g}$ iniciales al $10\,\%$, alcanzar $15\,\%$ exige añadir $29,412\,\mathrm{g}$
de agua, conservando la masa seca. Faltan ensayos para validar presión, torque y calidad del pellet.
La trazabilidad, sensibilidades y verificaciones nominales se documentan en el anexo independiente.
